\section{Independent certificates for optimality and infeasibility}
\label{sec:certificates}

The optimizer need not be trusted to have generated a complete kernel basis.
An exact optimum has a small certificate checked by elementary modular
identities.  Let $\mu\mid m$ be the proposed optimum and put $s=m/\mu$.
Supply $z_*\in R^n$ and a matrix $Y\in R^{\beta\times q}$ satisfying
\begin{align}
 Az_*&=b, & h_*&=c+Hz_*, &\gcd(m,h_{*,1},\ldots,h_{*,\beta})&=\mu,
 \label{eq:affine-primal-certificate}\\
 YA&=sH, &Yb+sc&=0\pmod m.
 \label{eq:affine-dual-certificate}
\end{align}

\begin{theorem}[Primal and dual optimality certificate]
Identities~\eqref{eq:affine-primal-certificate}--\eqref{eq:affine-dual-certificate}
certify that $S$ is nonempty and its exact minimum component count is $\mu$.
Every feasible input admits such a certificate for its true optimum.
An infeasible input admits a vector $y\in R^q$ satisfying
\begin{equation}
 yA=0,\qquad yb\ne0\pmod m.
 \label{eq:affine-infeasibility-certificate}
\end{equation}
All these certificates can be generated in deterministic polynomial time
without factoring $m$.
\end{theorem}

\begin{proof}
For every feasible $z$, the dual identities give
$s(c+Hz)=sc+YAz=sc+Yb=0\pmod m$.  Since $s=m/\mu$, each entry of
$c+Hz$ is divisible by $\mu$.  Every component count is thus a multiple
of $\mu$, and the primal witness attains $\mu$.  Conversely, use the
elementary annihilator identity
\begin{equation}
 \{\ell\in R^{1\times n}:\ell x=0\text{ for all }x\in\ker A\}
 =\{yA:y\in R^{1\times q}\}.
 \label{eq:finite-annihilator}
\end{equation}
Here is a proof valid also when $m$ is composite.  Write a Smith
decomposition $PAQ=D$ with $P,Q$ unimodular over $\mathbb Z$, hence
invertible modulo $m$.  For a nonzero diagonal entry $d_i$, put
$g_i=\gcd(m,d_i)$.  The corresponding kernel coordinate consists of
multiples of $m/g_i$, so an annihilating row must have its $i$th entry
divisible by $g_i$.  That is exactly the ideal generated by $d_i$ in
$R$.  A zero diagonal coordinate permits every kernel value and forces
the row entry to be zero.  Transforming back proves
\eqref{eq:finite-annihilator}.

For the true optimum $\mu$, Theorem~\ref{thm:attained-affine-content}
implies that every feasible holonomy vector has all entries divisible by
$\mu$.  Each row of $sH$ therefore annihilates $\ker A$, so
\eqref{eq:finite-annihilator} supplies the corresponding row of $Y$.
Evaluating at any feasible $z_0$ also gives $Yb+sc=0$.
The rows of $Y$ are found constructively by solving
$A^Ty_j^T=sH_j^T$ with the incremental congruence solver.  When $\mu=1$,
one may simply take $Y=0$.

For infeasibility, apply~\eqref{eq:finite-annihilator} to $A^T$:
if $b$ annihilated all of $\ker A^T$, it would lie in $\operatorname{im} A$,
contrary to infeasibility.  Solve the homogeneous system $A^Ty^T=0$,
then scan its generator columns for one whose dot product with $b$ is
nonzero.  Such a column must occur.  Equation~\eqref{eq:affine-infeasibility-certificate}
contradicts $Az=b$ directly.  Every construction used here is polynomial.
\end{proof}

The certificate verifier uses only modular matrix products and gcds; it
does not invoke the congruence solver or the optimizer.  The separately
reported number of feasible assignments is exact by the preceding
proposition, but is not part of this minimal certificate's contract.

